\chapter*{Preface}
\addcontentsline{toc}{chapter}{Preface}
This book presents practical methods for modelling and analysing electrical power systems with LTspice. The emphasis is placed on reproducible examples that connect physical interpretation, manual calculation and numerical simulation.

The material is intended for students, engineers and technically interested readers. Each chapter combines a concise theoretical introduction with worked calculations, circuit models, simulation results and practical remarks.

The book presents nine simulation examples, ranging from simple to more complex cases. The formatting of the book and, in some cases, the translation from German to English were done with the help of ChatGPT. All examples and technical details are, of course, provided by the author. I wish readers every success with their simulations and hope that they, too—like me—will find great enjoyment in them.
I would also like to thank the developers of LTspice for this excellent software, which made it possible for me to write this book. Although it is designed for simulating electronic circuits, it is so flexible that it allowed me to perform high-voltage analyses with it as well. Using a mathematical trick, I was also able to adapt LTspice for steady-state load-flow calculations, which led to the development of my second toolbox. This toolbox is also freely available and will be presented in more detail in a future book.

LTspice is a simulator developed by Analog Devices and is available free of charge. Having worked for more than 35 years as a load dispatcher in Vienna’s electricity distribution system, I have drawn on my practical experience in developing these toolboxes.
I hope that you enjoy exploring electrical power systems through simulation as much as I do. Some simulations can be challenging, but the results offer valuable insights. Good luck with your simulations!

\vspace{2cm}
\begin{flushright}
Friedrich Schmidt\\
Austria, July 2026
\end{flushright}
%\cleardoublepage
\newpage